\documentclass[11pt]{article}

\usepackage[margin=1in]{geometry}
\usepackage{amsmath,amssymb,amsthm}
\usepackage{graphicx}
\usepackage[hidelinks]{hyperref}

\newtheorem{theorem}{Theorem}
\newtheorem{lemma}{Lemma}
\newtheorem{remark}{Remark}

\newcommand{\N}{\mathbb N}
\newcommand{\elltwo}{\ell^2}

\title{A Literal Counterexample to ``Non-Weak-Embedding Implies Expander''}
\author{Run fa\_banach\_001, lane 7}
\date{2026-06-20}

\begin{document}
\maketitle

\section*{Status}

This packet gives a likely valid counterexample to the literal wording of the
last open question in Arzhantseva--Tessera~\cite{AT}: a bounded-degree sequence
of finite connected graphs need not be an expander merely because it does not
weakly embed into \(\ell^2\).

The result should be read as a wording/normalization counterexample.  It does
not address the more interesting strengthened question in which one rules out
linear-size expander pieces or asks for a structural characterization modulo
such pieces.

\section*{Source Problem Evidence}

The source question appears in Section 8, ``Open problems'', on page 20 of the
arXiv PDF:

\begin{center}
\includegraphics[width=0.95\linewidth]{figures/open_problem_crop.png}
\end{center}

\section*{Transcription}

After observing that expanders, and even large subsets of expanders, do not
weakly embed into \(\ell^p\), while the paper's relative expander examples do
weakly embed into \(\ell^2\), the authors ask whether a sequence of finite
graphs with uniformly bounded degree that does not weakly embed into
\(\ell^2\) is necessarily an expander.

\section*{Definitions}

A sequence of maps \(f_n:X_n\to \ell^2\) is a weak embedding if there is a
constant \(K\) such that every \(f_n\) is \(K\)-Lipschitz and, for every fixed
\(R>0\),
\[
  \lim_{n\to\infty}
  \sup_{x\in X_n}
  \frac{|f_n^{-1}(B(f_n(x),R))|}{|X_n|}
  =0.
\]
An expander sequence is a sequence of finite connected graphs of uniformly
bounded degree, with \(|X_n|\to\infty\), whose Cheeger constants are bounded
below by a positive constant.

\section*{Proof Intuition}

The obstruction to weak embedding is inherited from any linear-size expander
subgraph.  If a graph \(X_n\) contains an isometric copy of an expander
\(E_n\) with \(|E_n|\ge c|X_n|\), then any uniformly Lipschitz map
\(X_n\to\ell^2\) collapses a positive proportion of \(E_n\), hence a positive
proportion of \(X_n\), into one bounded Hilbert ball.  But this has nothing to
do with the global Cheeger constant of \(X_n\).  Attaching a long path to
\(E_n\) by one edge destroys expansion while preserving the linear-size
expander obstruction to weak embedding.

\begin{lemma}\label{lem:collapse}
Let \((E_n)\) be a bounded-degree expander sequence.  For every \(K>0\) there
are \(R>0\) and \(c>0\) such that, for every sequence of \(K\)-Lipschitz maps
\(f_n:E_n\to\ell^2\), for all \(n\) sufficiently large there is a vertex
\(v_n\in E_n\) with
\[
  |f_n^{-1}(B(f_n(v_n),R))|\ge c|E_n|.
\]
\end{lemma}

\begin{proof}
The usual expander Poincare inequality gives a constant \(B\), depending only
on the expander family and its degree bound, such that every \(K\)-Lipschitz
map \(f:E_n\to\ell^2\) satisfies
\[
  \frac{1}{|E_n|^2}\sum_{x,y\in E_n}\|f(x)-f(y)\|^2 \le B K^2 .
\]
This is precisely the \(p=2\) form of the expander inequality recalled in
Section 2 of~\cite{AT}.  Therefore some vertex \(v\in E_n\) satisfies
\[
  \frac{1}{|E_n|}\sum_{y\in E_n}\|f(v)-f(y)\|^2 \le B K^2 .
\]
By Markov's inequality, at least half of the vertices \(y\in E_n\) obey
\(\|f(v)-f(y)\|\le \sqrt{2B}\,K\).  Take \(R=\sqrt{2B}\,K\) and \(c=1/2\).
\end{proof}

\begin{theorem}
There exists a sequence \((X_n)\) of finite connected graphs with uniformly
bounded degree such that:
\begin{enumerate}
  \item \((X_n)\) is not an expander sequence;
  \item \((X_n)\) does not weakly embed into \(\ell^2\).
\end{enumerate}
\end{theorem}

\begin{proof}
Let \((E_n)\) be any bounded-degree expander sequence, and write
\(m_n=|E_n|\).  Let \(P_n\) be a path with \(m_n\) vertices.  Choose a vertex
\(a_n\in E_n\) and an endpoint \(b_n\in P_n\), and form \(X_n\) by joining
\(a_n\) and \(b_n\) with one new edge.

The degrees of \(X_n\) are uniformly bounded.  Also \(X_n\) is connected and
\(|X_n|=2m_n\to\infty\).

The sequence \((X_n)\) is not an expander.  Indeed, let \(A_n\) be the terminal
half of the path \(P_n\), away from the attachment point.  Then
\(|A_n|\le |X_n|/2\), while only one edge leaves \(A_n\).  Hence
\[
  h(X_n)\le \frac{1}{|A_n|}\to 0.
\]

It remains to prove that \((X_n)\) does not weakly embed into \(\ell^2\).
The copy of \(E_n\) inside \(X_n\) is isometric: the attached path creates no
shortcut between two vertices of \(E_n\).  Suppose \(f_n:X_n\to\ell^2\) is a
sequence of \(K\)-Lipschitz maps.  Restrict \(f_n\) to \(E_n\).  By
Lemma~\ref{lem:collapse}, there are \(R>0\), \(c>0\), and vertices
\(v_n\in E_n\) such that, for infinitely many \(n\) and in fact eventually,
\[
  |f_n^{-1}(B(f_n(v_n),R))\cap E_n|\ge c|E_n|.
\]
Consequently
\[
  \sup_{x\in X_n}
  \frac{|f_n^{-1}(B(f_n(x),R))|}{|X_n|}
  \ge
  \frac{c|E_n|}{2|E_n|}
  =
  \frac{c}{2}
\]
along the same tail.  Thus the weak-embedding condition fails for every
uniform Lipschitz constant \(K\).  Therefore \((X_n)\) does not weakly embed
into \(\ell^2\).
\end{proof}

\section*{Verification Notes}

No computation is needed.  The proof reduces the non-weak-embedding assertion
to the standard expander Poincare inequality and uses an explicit one-edge
bottleneck to force \(h(X_n)\to 0\).  The construction also shows a slightly
more general principle: any bounded-degree graph sequence containing a
linear-size isometric expander piece need not weakly embed into \(\ell^2\),
regardless of its global Cheeger constant.

\section*{Limitations}

This packet answers only the literal formulation displayed above.  It is not a
new relative-expander construction, and it does not address the intended
frontier suggested by~\cite{AT}: understanding non-weak-embeddability once
linear-size expander pieces have been excluded.  A natural repaired version
would ask whether a bounded-degree sequence with no linear-size expander
pieces, or no large subsets satisfying the expander Poincare collapse, must
weakly embed into \(\ell^2\).

\section*{Novelty and Literature Check}

Before promotion, I searched the run indexes for arXiv:1402.1481, ``relative
expanders'', ``weakly embed into \(\ell^2\)'', and ``necessarily an
expander''.  No existing packet or attempt covered this literal final
question.  Web searches for the exact wording found the source paper and later
papers answering other open problems from Section 8, including the finitely
generated group problem~\cite{AT2} and the coarse Baum--Connes problem, but no
exact prior answer to this literal graph-sequence question.

The novelty confidence is deliberately modest.  The construction is elementary
and the source paper itself points out that large subsets of expanders obstruct
weak embedding.  The value of the packet is to record a precise counterexample
to the literal wording and to clarify the hypothesis needed for a more
interesting version.

\section*{Human Review Recommendation}

Check that the weak-embedding normalization in~\cite{AT} uses \(|X_n|\) in the
denominator, and that \(|E_n|/|X_n|=1/2\).  Then verify that the attached path
does not create shortcuts inside \(E_n\) and that the terminal half of the path
forces \(h(X_n)\to 0\).

\begin{thebibliography}{9}

\bibitem{AT}
G. Arzhantseva and R. Tessera,
\emph{Relative expanders},
arXiv:1402.1481, 2014.
\url{https://arxiv.org/abs/1402.1481}

\bibitem{AT2}
G. Arzhantseva and R. Tessera,
\emph{Admitting a coarse embedding is not preserved under group extensions},
arXiv:1605.01192, 2016.
\url{https://arxiv.org/abs/1605.01192}

\end{thebibliography}

\end{document}
